\clearpage\fancyhead[L]{\small GAUSS AN BODE.}\hypertarget{native-392-lower}{}\label{native-392-lower}
\begin{center}\textbf{GAUSS AN BODE.}\par\vspace{4pt}\hrule\vspace{4pt}\fontsize{10}{12}\selectfont Astronomisches Jahrbuch für 1818. Seite 229..232. Berlin 1815.\par\vspace{4pt}\hrule\vspace{4pt}\end{center}

\begin{flushright}
Göttingen, den 9. August 1815.
\end{flushright}

\par\indent Gerne teile ich Ihnen hier alle meine Kometenbeobachtungen mit. Das äußerst ungünstige Wetter und der Mondschein werden mich, fürchte ich, wohl daran hindern, weitere Beobachtungen hinzuzufügen, obwohl der Komet sonst noch einige Wochen gut zu beobachten wäre.

\SourceLeaf{393}\begin{center}\fontsize{10}{12}\selectfont\setlength{\tabcolsep}{5pt}\renewcommand{\arraystretch}{1.15}\begin{tabular}{l|c|c|c}
1815 & \shortstack{Mittlere Zeit\\in Göttingen} & \shortstack{Scheinbare\\Rektaszension} & \shortstack{Scheinbare nördliche\\Deklination}\\\hline
März 20 & $10^{\mathrm u}\;33^{\mathrm m}\;6^{\mathrm s}$ & $54^\circ\;7'\;1''$ & $39^\circ\;7'\;47''$\\
21 & $10\;11\;37$ & $54\;33\;21$ & $39\;36\;57$\\
25 & $9\;46\;38$ & $56\;28\;50$ & $41\;38\;5$\\
30 & $9\;50\;57$ & $59\;13\;2$ & $44\;10\;27$\\
April 2 & $9\;11\;55$ & $61\;2\;27$ & $45\;39\;57$\\
Juni 12 & $10\;53\;31$ & $165\;3\;3$ & $52\;8\;25$\\
30 & $10\;49\;12$ & $184\;39\;24$ & $39\;22\;5$\\
Juli 13 & $10\;57\;20$ & $194\;32\;52$ & $29\;53\;54$\\
27 & $10\;4\;59$ & $203\;0\;56$ & $20\;37\quad:$\\
29 & $10\;13\;58$ & $204\;6\;3$ & $19\;24\;13$\\
Aug. 4 & $10\;14\;7$ & $207\;14\;8$ & $15\;53\;49$\\
\end{tabular}\end{center}

\par\indent Am 13. und 27. Juli sowie am 4. August wurden Sterne aus \textsc{Piazzis} Katalog (XII, 268; 1 Boot.; XIII, 264. 265. 284) zum Vergleich herangezogen, an allen anderen Tagen Sterne aus der Histoire Céleste und den Mem. 1790. Alle Beobachtungen wurden mit dem Kreismikrometer des 10 Fuß langen \textsc{Herschel}schen Teleskops angestellt, ausgenommen die vom 12. Juni. An diesem Tag wurde der Komet mit dem \textsc{Frauenhofer}schen Heliometer, das auf einer ausgezeichneten parallaktischen Montierung stand, mit 219 Ursae majoris \textsc{Bode} (berechnet aus den Mem. 1790) verglichen. Dazu dienten der Abstand und der unmittelbar am Instrument abgelesene Positionswinkel. So ausgezeichnet die Ergebnisse dieser Methode bei helleren Kometen auch sein mögen, in diesem Fall ziehe ich gute Beobachtungen mit dem Kreismikrometer vor. Der lichtschwache Komet verschwand neben dem mehr als 5-mal so hellen Stern völlig und war nur zu erkennen, wenn die eine Hälfte des Objektivs größtenteils abgedeckt war.

\Needspace{23\baselineskip}\par\indent Meine schon Ende März berechneten parabolischen Elemente habe ich Ihnen, glaube ich, bereits mitgeteilt. Sie weichen in den folgenden Monaten stärker von der Bewegung des Kometen ab, als man es sonst in solchen Fällen gewohnt ist. Sobald ich die Beobachtung vom 12. Juni erhalten hatte, berechnete ich aus ihr und zwei Beobachtungen von \textsc{Olbers} vom 6. März und 25. April neue parabolische Elemente, die ich hier angebe:

\Needspace{12\baselineskip}\begin{center}\fontsize{10}{12}\selectfont Durchgang durch das Perihel am 25. April 1815.\par $11^{\mathrm u}\;41^{\mathrm m}\;19^{\mathrm s}$ M. Z. in Göttingen\end{center}
\begin{center}\fontsize{9}{11}\selectfont\setlength{\tabcolsep}{2pt}\renewcommand{\arraystretch}{1.15}\begin{tabular}{lcr}
Länge des Perihels & $.\;.\;.\;.\;.\;.\;.\;.\;.\;.\;.\;.$ & $147^\circ\;35'\;55''$\\
Länge des aufsteigenden Knotens & $.\;.\;.\;.\;.\;.\;.\;.$ & $82\;43\;6$\\
Neigung der Bahn & $.\;.\;.\;.\;.\;.\;.\;.\;.\;.\;.\;.\;.$ & $44\;43\;13$\\
Kleinster Abstand von der Sonne & $.\;.\;.\;.\;.\;.\;.\;.$ & $1.23024$\\
\end{tabular}\end{center}\begin{center}Bewegung rechtläufig.\end{center}

\par\indent Diese Elemnnte geben die drei zugrunde gelegten Beobachtungen vollständig und fast genau wieder. Trotzdem weichen sie von den dazwischenliegenden Beobachtungen beträchtlich ab (in der Rektaszension um bis zu 11 Minuten). Es kann oft vorkommen, dass drei vollständige Beobachtungen nicht ausreichen, um den Kegelschnitt zu bestimmen. Darüber habe ich in den letzten Artikeln des ersten Abschnitts des zweiten Buchs der Theoria Motus C. C. ausführlicher gesprochen. Da aber die Abweichung der parabolischen Elemente von den anderen Beobachtungen eindeutig zeigte, dass sich die Bahn stark von einer Parabel unterscheiden musste, schien es mir interessant, eine vorläufige hypothesenfreie Berechnung der Bahn vorzunehmen. Ich legte 5 Beobachtungen zugrunde, die durch die daraus ermittelte Ellipse bis auf Kleinigkeiten wiedergegeben wurden. Auch haben diese Elemente (um die Mitte Juni berechnet und im 105. Stück der hiesigen Gel. Anz. [S. 385 d. B.] abgedruckt) bis auf meine letzten Beobachtungen durchgehend gut mit den Beobachtungen \SourceLeaf{394}übereingestimmt. Herr \textsc{Nicolai} hat sie inzwischen schon einmal durch eine genauere Rechnung und die Verwendung weiterer Beobachtungen verbessert. Er erhielt nur kleine Änderungen. Mit Freude sehe ich auch, dass unser ausgezeichneter \textsc{Bessel} Ergebnisse erhalten hat, die sehr gut damit übereinstimmen. Den weiteren Feinschliff werde ich diesen geschickten Händen überlassen. Die Ergebnisse werden ein kostbares Vermächtnis für unsere Enkel sein, die die Wiederkehr dieses höchst bemerkenswerten Kometen feiern werden.

\Needspace{14\baselineskip}\begin{center}
\emph{Elliptische Elemente des Kometen von} 1815.
\end{center}

\begin{center}\fontsize{10}{12}\selectfont Durchgang durch das Perihel am 26. April.\par $0^{\mathrm u}\;21^{\mathrm m}\;25^{\mathrm s}$ in Göttingen.\end{center}
\begin{center}\fontsize{9}{11}\selectfont\setlength{\tabcolsep}{2pt}\renewcommand{\arraystretch}{1.15}\begin{tabular}{lcr}
Länge des Perihels & $.\;.\;.\;.\;.\;.\;.\;.\;.\;.\;.\;.$ & $148^\circ\;58'\;48''$\\
Abstand vom Perihel & $.\;.\;.\;.\;.\;.\;.\;.\;.$ & $1.21349$\\
Aufsteigender Knoten & $.\;.\;.\;.\;.\;.\;.\;.\;.\;.\;.\;.$ & $83^\circ\;26'\;21''$\\
Neigung der Bahn & $.\;.\;.\;.\;.\;.\;.\;.\;.\;.\;.\;.\;.$ & $44\;30\;43$\\
Exzentrizität & $.\;.\;.\;.\;.\;.\;.\;.\;.\;.\;.\;.\;.\;.\;.$ & $0.933149$\\
\end{tabular}\end{center}

\par\indent Daraus ergeben sich ferner ein Abstand im Aphel von $35.0911$ und eine Umlaufzeit von $77\tfrac12$ Jahren. Diese Umlaufzeit kann nur um wenige Jahre unsicher sein und ist wahrscheinlich noch etwas zu groß.

\Needspace{9\baselineskip}\par\indent Von Juno, die bei der diesjährigen Opposition nur die 10. Größe hatte, habe ich die folgenden drei Beobachtungen gemacht. Vielleicht sind es überhaupt die einzigen, die diesmal angestellt wurden.

\begin{center}\fontsize{9}{11}\selectfont\setlength{\tabcolsep}{2pt}\renewcommand{\arraystretch}{1.15}\begin{tabular}{llrccc}
1815 & März & 1. & $10^{\mathrm u}\;22^{\mathrm m}\;0^{\mathrm s}$ & $\mathrm{AR}=197^\circ\;7'\;56''8$ & $\mathrm{Decl.}=2^\circ\;19'\;3''0\;\mathrm{S.}$\\
 & & 29. & $10\;8\;57$ & $192\;31\;51.8$ & $1\;32\;38.1\;\mathrm{N.}$\\
 & April & 8. & $9\;27\;43$ & $190\;34\;51.9$ & $2\;53\;24.2$ ---\\
\end{tabular}\end{center}

\par\indent Was das \textsc{Frauenhofer}sche Heliometer betrifft, für das ich im vergangenen Frühjahr das überaus schöne parallaktische Stativ erhalten habe, so wäre eine etwas ausführlichere Beschreibung ohne genaue Zeichnung wohl kaum verständlich. Ich bin überzeugt, dass man sich bei einer sehr festen Aufstellung des Instruments auf die Bewegungen der parallaktischen Montierung bis auf $\tfrac12$ Bogenminute verlassen könnte (die Nonien zeigen sogar $10''$ an). Das erlauben meine jetzigen Räumlichkeiten aber nicht. Ich muss das Instrument bald in diese, bald in jene Türöffnung stellen. Inzwischen habe ich jedoch einen Weg gefunden, es überall in sehr kurzer Zeit bis auf 1 oder höchstens $2'$ richtig auszurichten. Bei günstiger Luft finde ich dann zum Beispiel den Polarstern, den Merkur und dergleichen auch am Tag ohne große Mühe.
